% Options for packages loaded elsewhere
\PassOptionsToPackage{unicode}{hyperref}
\PassOptionsToPackage{hyphens}{url}
\documentclass[
]{article}
\usepackage{xcolor}
\usepackage{amsmath,amssymb}
\setcounter{secnumdepth}{-\maxdimen} % remove section numbering
\usepackage{iftex}
\ifPDFTeX
  \usepackage[T1]{fontenc}
  \usepackage[utf8]{inputenc}
  \usepackage{textcomp} % provide euro and other symbols
\else % if luatex or xetex
  \usepackage{unicode-math} % this also loads fontspec
  \defaultfontfeatures{Scale=MatchLowercase}
  \defaultfontfeatures[\rmfamily]{Ligatures=TeX,Scale=1}
\fi
\usepackage{lmodern}
\ifPDFTeX\else
  % xetex/luatex font selection
\fi
% Use upquote if available, for straight quotes in verbatim environments
\IfFileExists{upquote.sty}{\usepackage{upquote}}{}
\IfFileExists{microtype.sty}{% use microtype if available
  \usepackage[]{microtype}
  \UseMicrotypeSet[protrusion]{basicmath} % disable protrusion for tt fonts
}{}
\makeatletter
\@ifundefined{KOMAClassName}{% if non-KOMA class
  \IfFileExists{parskip.sty}{%
    \usepackage{parskip}
  }{% else
    \setlength{\parindent}{0pt}
    \setlength{\parskip}{6pt plus 2pt minus 1pt}}
}{% if KOMA class
  \KOMAoptions{parskip=half}}
\makeatother
\setlength{\emergencystretch}{3em} % prevent overfull lines
\providecommand{\tightlist}{%
  \setlength{\itemsep}{0pt}\setlength{\parskip}{0pt}}
\usepackage{bookmark}
\IfFileExists{xurl.sty}{\usepackage{xurl}}{} % add URL line breaks if available
\urlstyle{same}
% fallback for those not using the hyperref driver hyperxmp:
\makeatletter
\@ifundefined{xmpquote}{\newcommand{\xmpquote}[1]{#1}}{}
\makeatother
\hypersetup{
  hidelinks,
  pdfcreator={LaTeX via pandoc}}

\author{}
\date{}

\begin{document}

{
\setcounter{tocdepth}{2}
\tableofcontents
}
\section{WebCoder --- Teacher Guide}\label{webcoder--teacher-guide}

\begin{quote}
One-paragraph summary: WebCoder gives teachers the pbinfo-style
classroom loop --- classes with student self-join, assignments (teme)
with publish windows and soft expiry, exam-mode lockdown, MCQ sets from
the Bac bank, per-student progress grids, heatmaps, effort boards,
report cards and CSV/PDF exports --- on top of a real bilingual judge.
This guide walks the whole classroom workflow.
\end{quote}

\textbf{Who it\textquotesingle s for:} teachers, club coordinators,
county olympiad trainers. Romanian twin:
\href{../guide-ro/GHID-PROFESOR.md}{\texttt{docs/guide-ro/GHID-PROFESOR.md}}.

\begin{center}\rule{0.5\linewidth}{0.5pt}\end{center}

\subsection{1. Get the TEACHER role}\label{1-get-the-teacher-role}

Ask an administrator to upgrade your account to \textbf{TEACHER}
(verification via school email / the school-partnership program;
\texttt{docs/RUNBOOK.md} §Wave C). The role unlocks \texttt{/teach},
class creation, assignment authoring and the analytics surfaces.

\subsection{2. Create a class}\label{2-create-a-class}

\begin{enumerate}
\def\labelenumi{\arabic{enumi}.}
\tightlist
\item
  \texttt{/teach} → \textbf{New class} → name it (e.g. "IX-A
  Informatică").
\item
  You get a \textbf{join code + join password}. Students join themselves
  from their profile (pbinfo\textquotesingle s anti-abuse model --- you
  never create student accounts):

  \begin{itemize}
  \tightlist
  \item
    share the code, students enter it on \texttt{/classroom} join;
  \item
    or share the \textbf{QR code} (mobile-friendly join);
  \item
    batch import by username list is also available (missing usernames
    reported).
  \end{itemize}
\item
  Members list links to public profiles and class-scoped submission
  views.
\end{enumerate}

\subsection{3. Assignments (teme)}\label{3-assignments-teme}

Create an assignment on a class:

\begin{itemize}
\tightlist
\item
  \textbf{problems} picked individually, from curriculum chapters, or
  from any track in one click ("turn this track into a tema");
\item
  \textbf{availability window} (\texttt{available\ from} →
  \texttt{due\ at}, optional \texttt{late\ until});
\item
  optional \textbf{exam mode}: while a student is inside an exam temă,
  the rest of the site (forum, other problems, hints) is locked down for
  them, with a published conduct/AI policy;
\item
  \textbf{MCQ sets} pulled from the Bacalaureat grile bank
  (discipline/class/Bac-year filters, shuffle option).
\end{itemize}

Soft expiry (the pbinfo semantic): after the deadline students can still
submit --- those submissions simply don\textquotesingle t count for the
temă. The judge never punishes learning; upsolving doubles as revision.

\subsection{4. Track the class}\label{4-track-the-class}

The class page\textquotesingle s tabs:

\begin{itemize}
\tightlist
\item
  \textbf{Grid} --- students × problems matrix of best score / verdict /
  attempts, aggregated over submissions inside the window; per-student
  drawer with their full history in the class.
\item
  \textbf{Insights} --- class failure heatmap (which problems block
  whom), per-problem failure rates.
\item
  \textbf{Effort board} --- attempts + active days + solves ranking (the
  grind board).
\item
  \textbf{Report card} --- per-student PDF: every temă\textquotesingle s
  per-problem scores, totals, late flags. Printable for the catalog
  electronic.
\item
  \textbf{Exam PDF} --- assignments export as a printable
  statements+points booklet.
\end{itemize}

\subsection{5. Exports \& grades}\label{5-exports--grades}

\begin{itemize}
\tightlist
\item
  CSV gradebook (per-student max score per problem + totals) for any
  external catalog;
\item
  per-class report-card PDFs;
\item
  "download all solutions" per temă (the \#1 teacher request on every
  judge) for offline plagiarism review (plus the
  platform\textquotesingle s own advisory similarity flags).
\end{itemize}

\subsection{6. Curriculum tree}\label{6-curriculum-tree}

\texttt{/problems} curriculum axis maps the national programme (grades
5--12 → chapters → topics) with per-node problem counts. Teme-building
is chapter-picking: filter the catalog by your chapter, select, assign.
The tree is versioned data (a new programme edition adds a taxonomy
version, not a migration).

\subsection{7. Run a class contest}\label{7-run-a-class-contest}

Windowed contests fit a school schedule: one multi-day window, each
student gets one timed sitting (personal countdown). Create it from
\texttt{/contests} (TEACHER/CREATOR can create private contests), attach
the class\textquotesingle s problems, share the registration. Rated or
unrated --- unrated is the default for practice rounds. The replay page
shows any student\textquotesingle s minute-by-minute timeline for
post-mortem teaching.

\subsection{8. Authoring your own
problems}\label{8-authoring-your-own-problems}

CREATOR/TEACHER accounts can author:

\begin{itemize}
\tightlist
\item
  the 7-step \textbf{visual wizard} (basics → limits → samples → tests →
  reference solution → translation → review) with the statement linter
  and readability score;
\item
  hidden tests via the generator DSL or the testcase bounty board
  (community pledges);
\item
  Polygon/Kattis package import (permission-ledger gated).
\end{itemize}

Every problem goes through DRAFT → PENDING → APPROVED with verifier
feedback; authoring IS professional development (the contributor program
page documents the specs and feedback templates).

\subsection{9. Certifications \&
partnerships}\label{9-certifications--partnerships}

\begin{itemize}
\tightlist
\item
  Teacher certification (the program page explains requirements and the
  exam);
\item
  school partnerships (DPA template in \texttt{docs/DPA-TEMPLATE.md});
\item
  data-processing: classroom data follows the GDPR section of the legal
  pages; minors\textquotesingle{} consent is school-mediated
  (teacher-created accounts under the school\textquotesingle s lawful
  basis) or parental.
\end{itemize}

\subsection{10. FAQ}\label{10-faq}

\textbf{A student forgot their password} --- self-service reset by
email; without email, a school-mediated account is the robust pattern.
\textbf{Can I see students\textquotesingle{} code?} Yes --- class-scoped
submission views per student (the gradebook links them). \textbf{Do
contest ratings affect class grades?} No --- ratings and teme scores are
separate systems; export what you need.

\end{document}
